\documentclass[../../../include/open-logic-section]{subfiles}

\begin{document}
\olfileid[id]{sth}{card-arithmetic}{expotough}

\olsection[Beberapa Penyederhanaan]{Beberapa Penyederhanaan dengan Eksponensiasi Kardinal}

Walaupun mendefinisikan $\canonord$ agak rumit, hasil akhirnya ialah suatu
pernyataan yang berguna mengenai penjumlahan dan perkalian kardinal,
\olref[simp]{cardplustimesmax}. Akan tetapi, eksponensiasi transfinit tidak
dapat disederhanakan dengan begitu langsung. Untuk menjelaskan alasannya,
kita mulai dengan suatu hasil yang memperluas pola yang dikenal dari kasus
hingga (walaupun buktinya berada pada tingkat abstraksi yang tinggi):

\begin{prop}\ollabel{simplecardexpo}
$\cardexpo{\cardfont{a}}{\cardfont{b} \cardplus \cardfont{c}} =
\cardexpo{\cardfont{a}}{\cardfont{b}} \cardtimes
\cardexpo{\cardfont{a}}{\cardfont{c}}$ dan
$\cardexpo{(\cardexpo{\cardfont{a}}{\cardfont{b}})}{\cardfont{c}} =
\cardexpo{\cardfont{a}}{\cardfont{b} \cardtimes \cardfont{c}}$, untuk
sebarang kardinal $\cardfont{a}, \cardfont{b}, \cardfont{c}$.
\end{prop}

\begin{proof}
Untuk klaim pertama, perhatikan suatu fungsi $f \colon
(\cardfont{b}\disjointsum\cardfont{c}) \to \cardfont{a}$. Sekarang
``pisahkan fungsi ini'' dengan mendefinisikan
$f_\cardfont{b}(\beta) = f(\beta, 0)$ untuk setiap
$\beta \in \cardfont{b}$, dan
$f_\cardfont{c}(\gamma) = f(\gamma, 1)$ untuk setiap
$\gamma \in \cardfont{c}$. Pemetaan
$f \mapsto \tuple{f_{\cardfont{b}}, f_\cardfont{c}}$ merupakan
!!a{bijection}
$\funfromto{\cardfont{b} \disjointsum \cardfont{c}}{\cardfont{a}} \to
(\funfromto{\cardfont{b}}{\cardfont{a}} \times
\funfromto{\cardfont{c}}{\cardfont{a}})$.

Untuk klaim kedua, perhatikan suatu fungsi
$f \colon \cardfont{c} \to
(\funfromto{\cardfont{b}}{\cardfont{a}})$; jadi, untuk setiap
$\gamma \in \cardfont{c}$ kita mempunyai suatu fungsi
$f(\gamma) \colon \cardfont{b} \to \cardfont{a}$. Sekarang definisikan
$f^*(\beta, \gamma) = (f(\gamma))(\beta)$ untuk setiap
$\tuple{\beta, \gamma} \in \cardfont{b} \times \cardfont{c}$.
Pemetaan $f \mapsto f^*$ merupakan !!a{bijection}
$\funfromto{\cardfont{c}}{(\funfromto{\cardfont{b}}{\cardfont{a}})}
\to \funfromto{\cardfont{b} \cardtimes \cardfont{c}}{\cardfont{a}}$.
\end{proof}

Sekarang, yang kita \emph{inginkan} ialah suatu cara mudah untuk menghitung
$\cardexpo{\cardfont{a}}{\cardfont{b}}$ ketika kita menangani kardinal tak
hingga. Berikut suatu langkah yang baik ke arah tersebut:

\begin{prop}\ollabel{cardexpo2reduct}
Jika $2 \leq \cardfont{a} \leq \cardfont{b}$ dan $\cardfont{b}$ tak hingga,
maka $\cardexpo{\cardfont{a}}{\cardfont{b}} =
\cardexpo{2}{\cardfont{b}}$
\end{prop}

\begin{proof}
\begin{align*}
	\cardexpo{2}{\cardfont{b}} &\leq 
	\cardexpo{\cardfont{a}}{\cardfont{b}}\text{, karena $2 \leq \cardfont{a}$}\\
	&\leq \cardexpo{(2^\cardfont{a})}{\cardfont{b}}
	\text{, berdasarkan \olref[opps]{lem:SizePowerset2Exp}}\\
	&= \cardexpo{2}{\cardfont{a} \cardtimes \cardfont{b}}
	\text{, berdasarkan \olref{simplecardexpo}} \\
	&= \cardexpo{2}{\cardfont{b}}
	\text{, berdasarkan \olref[simp]{cardplustimesmax}}
\end{align*}
\end{proof}

Kita seharusnya tidak mengharapkan penyederhanaan lebih jauh, sebab
$\cardfont{b} < \cardexpo{2}{\cardfont{b}}$ berdasarkan
\olref[card-arithmetic][opps]{lem:SizePowerset2Exp}. Namun, hasil ini tidak
memberi tahu kita apa yang harus dikatakan tentang
$\cardexpo{\cardfont{a}}{\cardfont{b}}$ ketika
$\cardfont{b} < \cardfont{a}$. Tentu saja, jika $\cardfont{b}$
\emph{hingga}, kita mengetahui apa yang harus dilakukan.

\begin{prop}
Jika $\cardfont{a}$ tak hingga dan $0 < n \in \omega$, maka
$\cardexpo{\cardfont{a}}{n} = \cardfont{a}$
\end{prop}

\begin{proof}
$\cardexpo{\cardfont{a}}{n} = \cardfont{a} \cardtimes \cardfont{a}
\cardtimes \ldots \cardtimes \cardfont{a} = \cardfont{a}$, berdasarkan
\olref[simp]{cardplustimesmax}.
\end{proof}
\noindent
Selain itu, dalam beberapa kasus lain, kita dapat mengendalikan ukuran
$\cardexpo{\cardfont{a}}{\cardfont{b}}$:

\begin{prop}
Jika $2 \leq \cardfont{b} < \cardfont{a} \leq
\cardexpo{2}{\cardfont{b}}$ dan $\cardfont{b}$ tak hingga, maka
$\cardexpo{\cardfont{a}}{\cardfont{b}} = \cardexpo{2}{\cardfont{b}}$
\end{prop}

\begin{proof}
$\cardexpo{2}{\cardfont{b}}\leq \cardexpo{\cardfont{a}}{\cardfont{b}}
\leq \cardexpo{(\cardexpo{2}{\cardfont{b}})}{\cardfont{b}} =
\cardexpo{2}{\cardfont{b}\cardtimes\cardfont{b}} =
\cardexpo{2}{\cardfont{b}}$, dengan penalaran seperti dalam
\olref{cardexpo2reduct}.
\end{proof}
\noindent
Namun, setelah titik ini, persoalannya menjadi jauh lebih halus.

\end{document}
